\section*{Hoofdstuk \ref{ch:g12:synthesis-strategy} --- Syntheseplanning en groene chemie}

\begin{solution}{exo:g12:synthesis-strategy:1}
Hydratatie: $46.0/(28.0 + 18.0) = \qty{100}{\percent}$. Vergisting:
$2 \times 46.0/180.0 = \qty{51.1}{\percent}$.
\end{solution}

\begin{solution}{exo:g12:synthesis-strategy:2}
Beide groepen, \ce{Z} op de \omterm{def:g11:functional-groups:amine}{amine} van glycine en de methylester op het
zuur van alanine, worden in stap 1 aangebracht en in stap 3 verwijderd.
\end{solution}

\begin{solution}{exo:g12:synthesis-strategy:3}
Principe 5: gebruik veiligere \omterm{def:g6:solutions-and-solubility:solute}{oplosmiddelen} en reactieomstandigheden.
\end{solution}

\begin{solution}{exo:g12:synthesis-strategy:4}
Een overmaat ethanol maakt $Q < K$, dus de reactie gaat verder in de
heengaande richting. Een \omterm{def:g12:catalysis:catalyst}{katalysator} versnelt beide richtingen evenveel:
dezelfde \omterm{def:g10:reaction-progress-table:system}{eindtoestand} wordt bereikt, alleen eerder.
\end{solution}

\begin{solution}{exo:g12:synthesis-strategy:5}
Ja: alleen het \omterm{def:g11:functional-groups:carbonyl}{keton} wordt gereduceerd, tot een secundaire alcoholgroep;
de \omterm{def:g11:functional-groups:carboxylic-acid}{ester} blijft onveranderd.
\end{solution}

\begin{solution}{exo:g12:synthesis-strategy:6}
Meer \omterm{def:g11:functional-groups:carboxylic-acid}{ester}: een overmaat van een \omterm{def:g7:chemical-reactions:reactant}{beginstof}, of het water weghalen
(waterafscheider) zodra het ontstaat. Sneller: \omterm{def:g10:synthesis-yield:reflux}{verwarmen onder
terugvloeiing}, een zure \omterm{def:g12:catalysis:catalyst}{katalysator} toevoegen.
\end{solution}

\begin{solution}{exo:g12:synthesis-strategy:7}
Met \ce{NaBH4}:
\[ \ce{CH3-CH(OH)-CH2-CH2-CO-O-CH3} . \]
Met \ce{LiAlH4}:
\[ \ce{CH3-CH(OH)-CH2-CH2-CH2OH} \]
(pentaan-1,4-diol), plus methanol uit de andere helft van de \omterm{def:g11:functional-groups:carboxylic-acid}{ester}.
\end{solution}

\begin{solution}{exo:g12:synthesis-strategy:8}
De additie (\qty{100}{\percent}), de verestering
(\qty{83}{\percent}) en de route in drie stappen naar ibuprofen
(\qty{77}{\percent}). $77.4/40.0 = 1.9$: bijna twee keer zo goed.
\end{solution}

\begin{solution}{exo:g12:synthesis-strategy:9}
$180.0/(138.0 + 102.0) = \qty{75.0}{\percent}$. Ethaanzuur:
$60.0/180.0 = \qty{0.333}{kg}$ per kilogram aspirine.
\end{solution}

\begin{solution}{exo:g12:synthesis-strategy:10}
Het verwarmen versnelt de reactie; de koeler laat de dampen terugstromen
in de kolf, zodat er geen \omterm{def:g7:chemical-reactions:reactant}{beginstof} of \omterm{def:g6:solutions-and-solubility:solute}{oplosmiddel} verloren gaat.
Verwarmen verbetert de snelheid, niet het eindrendement.
\end{solution}

\begin{solution}{exo:g12:synthesis-strategy:11}
Zes stappen tegen drie. Principes: voorkom afval (1), maximaliseer de
\omterm{def:g12:synthesis-strategy:atom-economy}{atoomeconomie} (2), gebruik \omterm{def:g12:catalysis:catalyst}{katalysatoren} (9); ook minder stappen, minder
energie (6).
\end{solution}

\begin{solution}{exo:g12:synthesis-strategy:12}
$0.80 \times 0.70 \times 0.90 = 0.504$: \qty{50.4}{\percent}.
Uitgangsstof: $0.10/0.504 = \qty{0.20}{mol}$.
\end{solution}

\begin{solution}{exo:g12:synthesis-strategy:13}
Bescherm de \omterm{def:g11:functional-groups:amine}{amine} van alanine (\ce{Z}) en het zuur van glycine
(methylester); vorm de amidebinding tussen het vrije zuur van alanine en
de vrije \omterm{def:g11:functional-groups:amine}{amine} van glycine; verwijder beide \omterm{def:g12:synthesis-strategy:protecting-group}{beschermende groepen}. Drie
fasen, dat zijn vijf reacties als je elke bescherming en ontscherming
meetelt.
\end{solution}

\begin{solution}{exo:g12:synthesis-strategy:14}
De \omterm{def:g12:catalysis:catalyst}{katalysator} past principe 9 toe. Maar \qty{250}{\celsius} gaat in
tegen principe 6, benzeen tegen principes 3 en 5, en een \omterm{def:g12:synthesis-strategy:atom-economy}{atoomeconomie} van
\qty{35}{\percent} tegen principe 2: $65/35 = \qty{1.9}{kg}$ afval per
kilogram product. Eén groen kenmerk maakt nog geen groen proces.
\end{solution}

\begin{solution}{exo:g12:synthesis-strategy:15}
% ledger: crit:CO2-T, crit:CO2-p
\qty{40}{\celsius} en \qty{100}{bar} liggen boven het kritieke punt van
koolstofdioxide, \qty{31}{\celsius} en \qty{73.8}{bar}. De druk is
honderd keer de atmosferische druk: het vat is van dik staal.
Koolstofdioxide is niet giftig en niet brandbaar, het wordt weer gas
zonder resten achter te laten, en het kan opnieuw gebruikt worden; er
ontsnapt geen gechloreerd \omterm{def:g6:solutions-and-solubility:solute}{oplosmiddel}.
\end{solution}

\begin{solution}{pb:g12:synthesis-strategy:1}
% ledger: ibu:bhc-ae-recovered
\textbf{1.} C: $10 + 4 + 1 = 15 = 13 + 2$; H: $14 + 6 + 2 = 22 = 18 + 4$;
O: $3 + 1 = 4 = 2 + 2$.

\textbf{2.} Ethaanzuur, \ce{C2H4O2}.

\textbf{3.} Stikstof, chloor en natrium: ibuprofen bevat alleen
koolstof, waterstof en zuurstof.

\textbf{4.} De nieuwere: principe 9, \omterm{def:g12:catalysis:catalyst}{katalysatoren} in plaats van
stoichiometrische \omterm{def:g7:identifying-substances:characteristic-test}{reagentia}.

\textbf{5.} $13 \times 12.0 + 18 \times 1.0 + 2 \times 16.0 =
\qty{206.0}{g/mol}$.

\textbf{6.} $134.0 + 102.0 + 122.5 + 68.0 + 19.0 + 33.0 + 36.0 =
\qty{514.5}{g/mol}$.

\textbf{7.} $206.0/514.5 = \qty{40.0}{\percent}$.

\textbf{8.} $134.0 + 102.0 + 2.0 + 28.0 = \qty{266.0}{g/mol}$.

\textbf{9.} $206.0/266.0 = \qty{77.4}{\percent}$.

\textbf{10.} $(206.0 + 60.0)/266.0 = \qty{100}{\percent}$ op papier; meer
dan \qty{99}{\percent} in de praktijk.

\textbf{11.} $514.5/206.0 = \qty{2.50}{t}$.

\textbf{12.} $2.50 - 1 = \qty{1.50}{t}$.

\textbf{13.} $266.0/206.0 = \qty{1.29}{t}$.

\textbf{14.} $60.0/206.0 = \qty{0.29}{t}$.

\textbf{15.} $1.50 - 0.29 = \qty{1.21}{t}$.

\textbf{16.} $0.90^6 = 0.53$: \qty{53}{\percent}.

\textbf{17.} $0.90^3 = 0.73$: \qty{73}{\percent}.

\textbf{18.} Elke stap gebruikt \omterm{def:g6:solutions-and-solubility:solute}{oplosmiddelen}, scheidingen en
zuiveringen, en verliest wat product: minder stappen, minder van zulke
verliezen.

\textbf{19.} Bijna niets: haar enige bijproduct wordt gebruikt.

\textbf{20.} De route in drie stappen vermijdt ongeveer \qty{1.5}{t}
afval per ton ibuprofen (\qty{1.2}{t} zelfs als haar ethaanzuur zou
worden weggegooid).
\end{solution}
